\documentclass[10pt,a4paper]{amsart}
\usepackage[margin=19mm]{geometry}
\usepackage{amsmath,amssymb,mathtools}
\usepackage[scr=rsfs,cal=euler]{mathalfa}
\usepackage{xfrac}
\usepackage[unicode,hidelinks,bookmarks=false]{hyperref}
\newcommand{\ie}{i.e.\ }
\newcommand{\cf}{cf.\ }
\newcommand{\resp}{respectively\ }
\newcommand{\Subsection}{\subsection}
\providecommand{\nobreakdash}{\nobreak\hbox{-}}
\setlength{\emergencystretch}{2em}
\multlinegap0pt
\usepackage{amssymb}
\newtheorem{lemma}{Lemma}
\title{The BSE Property for Weighted Banach Algebras of Dirichlet Series}
\author{Prakash A. Dabhi}
\date{Source: arXiv:2609.15332v1 (CC BY 4.0)}
\begin{document}
\maketitle

\noindent\emph{Source and licence.} The BSE Property for Weighted Banach Algebras of Dirichlet Series. Version arXiv:2609.15332v1 (CC BY 4.0), distributed by arXiv under the Creative Commons Attribution 4.0 International licence, http://creativecommons.org/licenses/by/4.0/, as recorded in the arXiv licence metadata for this exact version and shown on the abstract page https://arxiv.org/abs/2609.15332v1. Full licence notice: \texttt{licenses/arxiv-2609.15332-CC-BY-4.0.txt}.\par\smallskip

\noindent\textit{Editorial scope.} Counted unit: the article's lemma lem:character\_approx (source0.tex lines 67--74). The article's complete original proof of it is delivered with it (source0.tex lines 76--86, one \texttt{proof} environment of 11 lines). No mathematics is written, repaired or re-derived: the statement and the proof are the article's own lines, in the article's own notation, and no retained line is substituted or re-flowed.  No further source-local context is needed: every cross-reference of every retained line resolves inside the two delivered spans.  Nothing was omitted from any retained span, so the package carries no omission record.  No retained line contains a citation, so the wrapper carries no bibliography and no bibliography entry.  No retained line contains a figure environment, an image-inclusion command or drawing code, so no image is delivered and the package declares no asset dependency.  Editorial formatting only, in addition: an AMS \texttt{amsart} preamble that restates the article's own declarations and macros used by the retained lines, namely the environments \texttt{lemma} on separate counters, together with the standard packages those lines need.  The retained environments print in this document as Lemma 1 in order of appearance, whereas the article prints the same items with the numbering of its own full document.  The complete native source of the article as distributed in the arXiv e-print for this exact version is bundled unchanged as \texttt{sources/210409/source0.tex}.
\begin{lemma}\label{lem:character_approx}
Let $\omega$ be a weight on $\mathbb N$. Fix $\varphi \in \Phi_{A_\omega}$, $k \in \mathbb N$ and $0 < r < 1$. Define the sequence $Z_{k,r} \in \mathbb C^{\mathbb N}$ by its values on primes as follows: $Z_{k,r}(p_i) = r \varphi(p_i)$ for $i \leq k$, and $Z_{k,r}(p_i) = 0$ for $i > k$. Extend it to all $n \in \mathbb N$ so that it is multiplicative. Define $\varphi_k \in \Phi_{A_\omega}$ to be the semicharacter matching $\varphi$ on the first $k$ primes and vanishing on all other primes. Then the following statements hold.
\begin{enumerate}
    \item The sequence $Z_{k,r}$ belongs to $\Phi_{A_\omega} \cap c_0(\mathbb N, 1/\omega)$.
    \item $Z_{k,r}\to\varphi_k$ in the weak-$\ast$ topology as $r \nearrow 1$.
    \item $\varphi_k\to\varphi$ in the weak-$\ast$ topology as $k \to \infty$.
\end{enumerate}
\end{lemma}

\begin{proof}
For the first claim, the multiplicative extension of $Z_{k,r}$ means that if $n=p_1^{\alpha_1}\cdots p_k^{\alpha_k}$ for some nonnegative integers $\alpha_1, \ldots, \alpha_k$, then we take $Z_{k,r}(n) = r^{\alpha_1+\cdots+\alpha_k}\varphi(n) = r^{\Omega(n)}\varphi(n)$ and if $p_j|n$ for some $j>k$, then $Z_{k,r}(n)=0$. Since $\varphi \in \Phi_{A_\omega}$, $|\varphi(n)| \leq \omega(n)$ for all $n\in \mathbb N$. As $r < 1$, we have $|Z_{k,r}(n)| \leq r^{\Omega(n)} \omega(n)$ for all $n \in \mathbb N$. As $n \to \infty$ within the support generated by the first $k$ primes, the exponent $\Omega(n) \to \infty$. As $0<r<1$, the ratio $Z_{k,r}(n) / \omega(n)\to 0$ as $n\to \infty$. It means that $Z_{k,r}$ is in $c_0(\mathbb N, 1/\omega)$. Because $|Z_{k,r}(n)| \leq \omega(n)$ for all $n$ and $Z_{k,r}$ is multiplicative, $Z_{k,r}$ is an $\omega$-bounded semicharacter, i.e., $Z_{k,r}\in \Phi_{A_\omega}$.

To prove the second claim, let $S_k$ denote the subsemigroup generated by the first $k$ primes. Let $g(s) = \sum_{n=1}^\infty a(n) n^{-s}$ be in $A_\omega$. Then
$$|Z_{k,r}(g) - \varphi_k(g)| \leq \sum_{n \in S_k} |a(n)| |\varphi(n)| (1 - r^{\Omega(n)}).$$
Because $|\varphi(n)| \leq \omega(n)$ for all $n$, the individual terms in this sum are bounded by $|a(n)| \omega(n)$. Since $g \in A_\omega$, $\sum_{n=1}^\infty |a(n)| \omega(n)<\infty$. For every fixed $n \in S_k$, the factor $(1 - r^{\Omega(n)})$ goes to zero as $r \nearrow 1$. It follows using the dominated convergence theorem that $\lim_{r\nearrow 1}\sum_{n \in S_k} |a(n)| |\varphi(n)| (1 - r^{\Omega(n)})=0$. This proves that $Z_{k,r}(g) \to \varphi_k(g)$ as $r\nearrow 1$. It means that $Z_{k,r}$ converges in the weak-$\ast$ topology to the semicharacter $\varphi_k$ as $r\nearrow 1$.

For the third claim, we take the limit as $k \to \infty$. For the same $g(s) \in A_\omega$, we have
$$|\varphi(g) - \varphi_k(g)| \leq \sum_{n \notin S_k} |a(n)| \omega(n).$$
Because $\|g\|_\omega=\sum_{n=1}^\infty|a(n)|\omega(n)<\infty$, the tail sum $\sum_{n \notin S_k} |a(n)| \omega(n) \to 0$ as $k \to \infty$. It means that $\varphi_k(g) \to \varphi(g)$ as $k\to \infty$. Therefore the sequence $(\varphi_k)$ converges to $\varphi$ in the weak-$\ast$ topology.
\end{proof}

\end{document}
